\chapter{Introducción}
\label{cap:introduccion}

\section{Motivación}

Los sistemas de conducción autónoma integran múltiples componentes de software que deben operar de manera coordinada para garantizar la seguridad. Entre estos componentes, los módulos de percepción basados en machine learning son particularmente difíciles de testear: no tienen especificaciones formales completas, su comportamiento se aprende de datos, y pueden fallar de formas sutiles e inesperadas.

El testing de estos sistemas presenta un desafío fundamental: ¿cómo diseñar tests efectivos para un componente cuyo comportamiento correcto no está formalmente especificado? Este trabajo propone que la respuesta está en el \textbf{conocimiento interno del sistema}. Si comprendemos profundamente cómo funciona un componente ---sus parámetros, sus reglas de decisión, sus casos límite--- podemos diseñar tests que exploten ese conocimiento para encontrar comportamientos anómalos.

Sin embargo, los sistemas de conducción autónoma como Apollo de Baidu son extremadamente complejos: millones de líneas de código, decenas de dependencias, configuraciones específicas de hardware. Esta complejidad dificulta tanto la comprensión del sistema como el diseño de tests controlados.

\section{Propuesta}

Este trabajo propone una metodología basada en \textbf{recortar y reimplementar} un subsistema crítico para poder comprenderlo profundamente y testearlo de forma controlada.

La idea central es:

\begin{enumerate}
    \item \textbf{Seleccionar un subsistema funcionalmente relevante} del sistema completo
    \item \textbf{Reimplementarlo} de forma fiel pero simplificada, eliminando complejidad accidental (dependencias, configuraciones)
    \item \textbf{Usar el conocimiento adquirido} durante la reimplementación para diseñar tests específicos
    \item \textbf{Buscar adversarios}: casos de test que revelen debilidades del sistema mediante análisis de caja blanca
\end{enumerate}

\subsection{¿Por qué reimplementar?}

El objetivo no es crear un sistema mejor, sino obtener un \textbf{subsistema equivalente sobre el cual tengamos control total}. Esto permite:

\begin{itemize}
    \item \textbf{Comprensión profunda:} Reimplementar fuerza a entender cada detalle del sistema original
    \item \textbf{Simplicidad:} Eliminar dependencias innecesarias facilita la experimentación
    \item \textbf{Instrumentación:} Agregar logging, visualizaciones, modificaciones para testing
    \item \textbf{Fidelidad:} Si la reimplementación es equivalente, los hallazgos aplican al sistema original
\end{itemize}

\subsection{¿Por qué buscar adversarios mediante caja blanca?}

La hipótesis central de este trabajo es que si tenemos acceso al código y comprendemos la lógica interna de un sistema, podemos diseñar tests más efectivos que con testing de caja negra. Específicamente:

\begin{itemize}
    \item Podemos identificar \textbf{parámetros críticos} (thresholds, pesos) y testear cerca de sus límites
    \item Podemos encontrar \textbf{reglas de negocio} y diseñar inputs que las violen
    \item Podemos descubrir \textbf{funcionalidades no implementadas} que podrían causar problemas
    \item Podemos entender \textbf{modos de falla} analizando la lógica de clasificación
\end{itemize}

\section{Caso de Estudio}

Como caso de estudio, seleccionamos el módulo de detección de semáforos (Traffic Light Recognition - TLR) del sistema Apollo de Baidu \cite{apollo_auto}.

\subsection{Criterios de Selección}

\begin{enumerate}
    \item \textbf{Criticidad:} La detección de semáforos es un componente de seguridad crítico

    \item \textbf{Simplicidad relativa:} Solo requiere cámaras (no LiDAR), tiene un espacio de salida limitado (rojo, amarillo, verde, apagado)

    \item \textbf{Interfaces claras:} Input (imagen + regiones de interés) → Output (detecciones + estados)

    \item \textbf{Código disponible:} Apollo es open source, permitiendo análisis completo
\end{enumerate}

\subsection{Apollo}

Apollo es una plataforma de código abierto para vehículos autónomos desarrollada por Baidu \cite{apollo_auto}. El módulo TLR es responsable de detectar semáforos en imágenes de cámara, reconocer su estado actual, y mantener tracking temporal de los semáforos a lo largo del tiempo.

\section{Contribuciones}

Las contribuciones de este trabajo son:

\begin{enumerate}
    \item \textbf{Metodología de recorte y reimplementación:} Demostración de que reimplementar un subsistema es más viable que intentar recortarlo del sistema complejo original

    \item \textbf{Análisis del sistema Apollo:} Documentación de decisiones de diseño no documentadas, parámetros extraídos, y hallazgos sobre el sistema

    \item \textbf{Criterios de conformidad:} Definición de qué significa que el recorte sea ``equivalente'' al original

    \item \textbf{Proceso de testing basado en conocimiento interno:} Diseño de tests específicos aprovechando la comprensión profunda del sistema

    \item \textbf{Validación de la hipótesis:} Demostración de que el análisis de caja blanca permite encontrar adversarios (caso del test de transición verde→amarillo)
\end{enumerate}

\section{Estructura del Trabajo}

Este trabajo se organiza de la siguiente manera:

\begin{itemize}
    \item \textbf{Capítulo \ref{cap:marco_teorico} - Marco Teórico:} Conceptos necesarios para comprender el trabajo, incluyendo todo lo conocido previamente sobre Apollo

    \item \textbf{Capítulo \ref{cap:analisis_apollo} - Análisis del Sistema:} El proceso de análisis de Apollo, las preguntas que guiaron el trabajo, y los hallazgos que justificaron las decisiones de recorte

    \item \textbf{Capítulo \ref{cap:reimplementacion} - Reimplementación:} Objetivos del recorte, alcance funcional, criterios de conformidad, y el proceso de reimplementación incluyendo la validación de equivalencia

    \item \textbf{Capítulo \ref{cap:diseno_tests} - Testing:} Diseño de tests basados en conocimiento interno del sistema, con foco en el test de transición verde→amarillo

    \item \textbf{Capítulo \ref{cap:resultados} - Resultados:} Hallazgos del proceso de testing y análisis

    \item \textbf{Capítulo \ref{cap:conclusiones} - Conclusiones:} Lecciones aprendidas, limitaciones, y trabajo futuro
\end{itemize}